\chapter{Parsing}
\label{ch:parsing}

Tokens tell the compiler which words appear, but not how they combine.
The text \code{2 + 3 * 4} and the text \code{(2 + 3) * 4} contain similar tokens yet denote different computations.
This chapter builds a \term{parser} that converts the token stream into the abstract syntax trees defined two chapters ago.
The parser enforces the grammar of Mini: which token orders form valid programs and how each valid order groups into a tree.
Everything after parsing works on trees and never re-examines grouping.

\section{Grouping decides meaning}
\label{sec:ch04-grouping-meaning}

Consider the expression \code{2 + 3 * 4}.
Two groupings are possible: \code{(2 + 3) * 4}, which equals $20$, and \code{2 + (3 * 4)}, which equals $14$.
Mini follows the ordinary mathematical convention that multiplication binds tighter than addition, so the correct tree is the second one.
A \term{grammar} writes such conventions as rules.
Each rule names a syntactic category and lists the token patterns that build it.
Mini expressions use three levels: sums, products, and atoms.
A sum is a product followed by any number of \code{+} or \code{-} steps, and a product is an atom followed by any number of \code{*} or division steps.
An atom is a number, a variable, or a parenthesized expression.
Because sums are built from products but not the reverse, multiplication always groups first.
Parentheses override the default by making a full expression act as one atom.
Table~\ref{tab:ch04-grammar} summarizes the syntactic categories; its verbal rows are an overview rather than a complete executable grammar.
\begin{table}
\caption{The Mini grammar with one rule per syntactic category and grouping fixed by rule layering.}
\label{tab:ch04-grammar}
\centering
\begin{tabular}{ll}
\toprule
Category & Forms \\
\midrule
Program & statement followed by more statements \\
Statement & assignment, let binding, print, conditional, loop, or block \\
Assignment & identifier, \code{=}, expression \\
Expression & product, then any number of add or subtract steps \\
Product & atom, then any number of multiply or divide steps \\
Atom & number, variable, or parenthesized expression \\
Conditional & \code{if} test \code{then} branch \code{else} branch \\
Loop & \code{while} test \code{do} body \\
\bottomrule
\end{tabular}
\end{table}
The grammar separates two questions that beginners often mix.
Whether a token order is valid is a question about the rules.
How a valid order groups is a question about the layering of expression and product.
The layered rules answer both at once: they accept exactly the valid orders and they admit exactly one grouping per order.
An \term{ambiguous} grammar would admit two trees for one token stream, and Mini avoids that by construction.
Associativity is the remaining choice: \code{10 - 3 - 2} groups as \code{(10 - 3) - 2}, which is $5$, because each layer collects repeated operators from left to right.
\section{Recursive descent}
\label{sec:ch04-recursive-descent}

A \term{recursive-descent parser} turns each grammar rule into one function.
The function for expressions calls the function for products, which calls the function for atoms, mirroring the grammar layering directly.
Each function consumes exactly the tokens of its category and returns the corresponding subtree.
A shared cursor holds the current token, and a helper advances it after each match.
When the next token cannot start the expected construct, the parser reports a syntax error naming the expected category and the token found.
Figure~\ref{fig:ch04-parse-tree} shows the tree that these functions build for \code{2 + 3 * 4}.
\begin{figure}
\centering
\begin{tikzpicture}[node distance=1.1cm]
\node[basic block] (plus) {Add};
\node[basic block, below left of=plus, xshift=-1.4cm] (two) {Num 2};
\node[basic block, below right of=plus, xshift=1.4cm] (mul) {Mul};
\node[basic block, below left of=mul, xshift=-1.0cm] (three) {Num 3};
\node[basic block, below right of=mul, xshift=1.0cm] (four) {Num 4};
\draw[cfg edge] (plus) -- (two);
\draw[cfg edge] (plus) -- (mul);
\draw[cfg edge] (mul) -- (three);
\draw[cfg edge] (mul) -- (four);
\end{tikzpicture}
\caption{The parse tree for \code{2 + 3 * 4} groups the multiplication first, so evaluation yields $14$.}
\label{fig:ch04-parse-tree}
\end{figure}
The root is the addition because sums sit at the top layer of the grammar.
Its right child is the multiplication, which sits one layer down and therefore binds tighter.
Evaluating this tree gives $3 \times 4 = 12$ and then $2 + 12 = 14$, matching the convention.
The alternative tree with multiplication at the root would evaluate to $20$ and corresponds to the parenthesized input \code{(2 + 3) * 4}.
Parentheses work by recursion: the atom function sees \code{(}, parses a full expression inside, and expects \code{)}.
The inner expression can contain any operators because it re-enters the top layer.
This single recursive call is the whole secret of nested grouping.

\section{The parser}
\label{sec:ch04-parser}

Algorithm~\ref{alg:ch04-parse} parses Mini expressions with correct precedence and left associativity.
\begin{algorithm}
\caption{Recursive-descent parsing of Mini expressions with layered precedence.}
\label{alg:ch04-parse}
\begin{pseudocode}
\Require token list $T$ with a cursor at position $P$
\Ensure syntax tree for one expression with the cursor past its last token, or a syntax error
\Function{ParseExpr}{$T$}
\State set $Left$ to the result of parsing one product
\While{the cursor shows an add or subtract operator}
\State consume the operator and parse the next product as $Right$
\State set $Left$ to a new node combining $Left$ and $Right$
\EndWhile
\State \Return $Left$
\EndFunction
\end{pseudocode}
\end{algorithm}
The algorithm assumes the cursor starts at the first token of an expression and that the product and atom parsers follow the same contract.
It guarantees that it consumes exactly the tokens of one expression: no fewer, so the tree is complete, and no more, so the caller keeps its place.
The loop collects repeated additions from left to right, which produces left-leaning trees for chains such as \code{10 - 3 - 2}.
Each iteration wraps the tree built so far as the left child of the new node, so earlier operators sit lower and evaluate first.
The product parser is a copy of this shape with multiply and divide operators calling the atom parser.
The displayed routine parses arithmetic sums.
To parse comparisons, wrap it with a routine that accepts at most one comparison operator and a second sum; this makes \code{x + 1 < y * 2} compare two arithmetic operands.
Parenthesized atoms call the full comparison parser, so \code{1 + (2 > 3)} parses and is rejected later by type checking.
The statement parser handles assignments, \code{let}, \code{print}, \code{if}, \code{while}, and blocks.
The caller must require the appropriate delimiter or end of input after a parsed expression; accepting a valid prefix does not validate the entire token stream.
Error reporting names the construct being parsed, the token found, and its position from the lexer.
Recovery in a batch compiler skips to the next statement boundary so that one mistake does not hide the rest.
An interactive mode may instead reparse the single line, since the cost of a retry is tiny.

\section{Worked example}
\label{sec:ch04-worked-example}

We now parse \code{2 + 3 * 4 - 1} step by step with Algorithm~\ref{alg:ch04-parse}.
The token list is \code{2}, \code{+}, \code{3}, \code{*}, \code{4}, \code{-}, \code{1}, which is $7$ tokens.
The cursor starts at position $0$.
Parsing the first product consumes only token $0$, because the next token is \code{+} rather than a multiplicative operator.
The value of $Left$ is the leaf \code{Num 2} and the cursor rests at position $1$.
The loop sees \code{+} and consumes it, moving the cursor to position $2$.
It parses the next product starting at token $2$: the atom \code{3} is followed by \code{*}, so the product loop consumes tokens $2$ through $4$ and builds \code{Mul(3, 4)}.
The expression loop wraps the tree into \code{Add(2, Mul(3, 4))} and the cursor rests at position $5$.
The loop sees \code{-} and consumes it, moving the cursor to position $6$.
It parses the final product, which is the single atom \code{1}, and wraps the tree into \code{Sub(Add(2, Mul(3, 4)), 1)}.
The loop then sees the end of input and returns.
The final tree has $3$ operator nodes and $4$ leaves, built in $7$ token advances and $7$ node constructions, of which $3$ construct operators.
Evaluating it gives $2 + 12 - 1 = 13$.
The count generalizes: parsing $K$ tokens performs exactly $K$ advances and builds at most $K$ nodes, which anticipates the linear bound of Section~\ref{sec:ch04-implementation}.
A second check confirms associativity: parsing \code{10 - 3 - 2} builds \code{Sub(Sub(10, 3), 2)}.
The inner subtraction evaluates to $7$ and the outer to $5$.
A right-leaning tree would instead give $10 - 1 = 9$, which Mini never builds for this input.

\section{Implementation notes}
\label{sec:ch04-implementation}

Each grammar rule becomes one function, so the parser stays proportional in size to Table~\ref{tab:ch04-grammar}.
The call stack mirrors the nesting of the input: deeply nested parentheses mean deep recursion.
Inputs generated by tools can nest thousands deep, so production parsers either raise the stack limit explicitly or rewrite nesting iteratively.
Left recursion needs care: a rule that calls itself before consuming any token would loop forever.
The layered grammar of this chapter already removes left recursion by using loops instead of self-calls for operator chains.
Error messages should quote the offending token and state what the parser expected at that point.
Messages that only say \enquote{syntax error} waste the positions the lexer worked to provide.
The running time is $O(K)$, where $K$ is the number of tokens.
Each token is examined once to decide the current step and consumed once when matched, and each loop iteration builds one tree node.
Memory use is $O(K + D)$, where $K$ bounds the size of the finished tree and $D$ is the maximum recursion depth reached while parsing nested constructs.
Generated parser tools trade the hand-written clarity of this section for automatic table construction from declarative grammars.
They handle more grammar shapes but produce harder-to-read errors unless the grammar author adds recovery rules.
For a language the size of Mini, the hand-written form of Algorithm~\ref{alg:ch04-parse} is shorter, clearer, and easier to extend with positions and recovery.
\begin{ccnote}[One tree per input]
A deterministic parser must build exactly one tree for each valid token stream.
If two groupings ever seem possible, the grammar is ambiguous and the fix belongs in the grammar, not in a tie-breaking hack inside the parser.
\end{ccnote}

\section{Exercises}
\label{sec:ch04-exercises}

\begin{enumerate}
\item \label{ex:ch04-trace-parse} Parse \code{x = 4 + 2 * 3;} with the rules of Table~\ref{tab:ch04-grammar}. Draw the resulting tree in the style of Figure~\ref{fig:ch04-parse-tree}, evaluate it with \code{x} unbound elsewhere, and state the final value. Confirm that the tree has $2$ operator nodes and $3$ leaves.
\item \label{ex:ch04-assoc} Parse \code{100 - 10 - 10 - 10} and show that Algorithm~\ref{alg:ch04-parse} builds a left-leaning tree evaluating to $70$. Write the right-leaning tree for the same tokens, evaluate it, and explain which grammar change would produce it.
\item \label{ex:ch04-error} Give three Mini token streams that are each invalid for a different reason: a missing operand, a missing terminator, and a misplaced keyword. For each one, state the exact position where parsing fails and the message your implementation should emit.
\end{enumerate}

\section{Further reading}
\label{sec:ch04-further-reading}

Appel develops parsing from grammars through recursive descent to table-driven methods, which places Algorithm~\ref{alg:ch04-parse} at the start of a well-marked road, see~\cite{appel-1998-modern}.
Morgan treats the front end as production engineering, where grammar clarity and error quality decide how fast users find their mistakes, see~\cite{morgan-1998-optimizing}.
The C language definition shows the same ideas at full scale, with a grammar whose precedence and associativity rules every C programmer relies on, see~\cite{wg14-2011-c11}.
Safonov connects front-end rigor to end-to-end trust in the compiler, which is why this chapter rejects ambiguity instead of papering over it, see~\cite{safonov-2010-trustworthy}.
